\documentclass[10pt,a4paper]{amsart}
\usepackage[margin=19mm]{geometry}
\usepackage{amsmath,amssymb,mathtools}
\usepackage[scr=rsfs,cal=euler]{mathalfa}
\usepackage{xfrac}
\usepackage[unicode,hidelinks,bookmarks=false]{hyperref}
\newcommand{\ie}{i.e.\ }
\newcommand{\cf}{cf.\ }
\newcommand{\resp}{respectively\ }
\newcommand{\Subsection}{\subsection}
\providecommand{\nobreakdash}{\nobreak\hbox{-}}
\setlength{\emergencystretch}{2em}
\multlinegap0pt
\usepackage[shortlabels]{enumitem}
\usepackage{amssymb}
\usepackage{xcolor}
\usepackage{float}
\usepackage{tikz}
\usepackage{algorithm}\usepackage{algpseudocode}
\newtheorem{restatable}{Restatable}
\title{Wasserstein Parallel Transport for Predicting the Dynamics of Statistical Systems}
\author{Tristan Luca Saidi, Gonzalo Mena, Larry Wasserman, Florian Gunsilius}
\date{Source: arXiv:2603.23736v1 (CC BY 4.0)}
\begin{document}
\maketitle

\noindent\emph{Source and licence.} Wasserstein Parallel Transport for Predicting the Dynamics of Statistical Systems. Version arXiv:2603.23736v1 (CC BY 4.0), distributed by arXiv under the Creative Commons Attribution 4.0 International licence, http://creativecommons.org/licenses/by/4.0/, as recorded in the arXiv licence metadata for this exact version and shown on the abstract page https://arxiv.org/abs/2603.23736v1. Full licence notice: \texttt{licenses/arxiv-2603.23736-CC-BY-4.0.txt}.\par\smallskip

\noindent\textit{Editorial scope.} Counted unit: the article's lemma lemma (source0.tex lines 1502--1555). The article's complete original proof of it is delivered with it (source0.tex lines 1557--1746, one \texttt{proof} environment of 190 lines). No mathematics is written, repaired or re-derived: the statement and the proof are the article's own lines, in the article's own notation, and no retained line is substituted or re-flowed.  No further source-local context is needed: every cross-reference of every retained line resolves inside the two delivered spans.  Nothing was omitted from any retained span, so the package carries no omission record.  No retained line contains a citation, so the wrapper carries no bibliography and no bibliography entry.  No retained line contains a figure environment, an image-inclusion command or drawing code, so no image is delivered and the package declares no asset dependency.  Editorial formatting only, in addition: an AMS \texttt{amsart} preamble that restates the article's own declarations and macros used by the retained lines, namely the environments \texttt{restatable} on separate counters, together with the standard packages those lines need.  The retained environments print in this document as Restatable 1 in order of appearance, whereas the article prints the same items with the numbering of its own full document.  The complete native source of the article as distributed in the arXiv e-print for this exact version is bundled unchanged as \texttt{sources/197103/source0.tex}.
\begin{restatable}[Time regularity of the weighted Helmholtz projection]{lemma}{}
\label{lem:time_reg_weighted_projection}
Let $\mu_t$ admit a Lebesgue density $\rho_t$ for $t \in [0, 1]$ such that:
\begin{enumerate}
    \item there exist constants
    $0<\lambda\le \Lambda<\infty$ such that for all $t\in[0,1]$,
    \[
    \lambda \le \rho_t(x)\le \Lambda
    \qquad\text{for a.e. }x\in\mathbb{T}^d.
    \]
    \item the map $t \mapsto \rho_t$ is absolutely continuous as a map from $[0,1]$ into $L^\infty(\mathbb{T}^d)$ and there exists
    $K_\rho<\infty$ such that,
    \[
    \sup_{t\in[0,1]}\|\partial_t \rho_t\|_{L^\infty(\mathbb{T}^d)}\le K_\rho.
    \]
\end{enumerate}
Then there exists a constant $C_{\dot \Pi}>0$, depending only on
$\lambda,\Lambda,$ and $K_\rho$, such that the following hold.
\begin{enumerate}
    \item For every $z\in L^2(\mathbb{T}^d;\mathbb R^d)$ and all
    $s,t\in[0,1]$,
    \[
    \|\Pi_{\mu_t}z-\Pi_{\mu_s}z\|_{L^2(\mathbb{T}^d)}
    \le
    C_{\dot\Pi}|t-s|\,\|z\|_{L^2(\mathbb{T}^d)}.
    \]
    In particular,
    \[
    \|\Pi_{\mu_t}-\Pi_{\mu_s}\|_{\mathcal L(L^2(\mathbb{T}^d;\mathbb R^d))}
    \le
    C_{\dot\Pi}|t-s|
    \]
    where the norm is the \textit{operator} norm on bounded linear maps $L^2(\mathbb{T}^d; \mathbb{R}^d) \rightarrow L^2(\mathbb{T}^d; \mathbb{R}^d).$
    \item For every absolutely continuous curve
    $z_\cdot \in AC([0,1];L^2(\mathbb{T}^d;\mathbb R^d))$, the curve
    $t\mapsto \Pi_{\mu_t}z_t$ belongs to
    $AC([0,1];L^2(\mathbb{T}^d;\mathbb R^d))$.
    Moreover, for a.e. $t\in[0,1]$, define
    \[
    (\partial_t\Pi_t)z_t
    \;\triangleq\;
    \partial_t(\Pi_{\mu_t}z_t)-\Pi_{\mu_t}(\partial_t z_t).
    \]
    Then
    \[
    \|(\partial_t\Pi_t)z_t\|_{L^2(\mathbb{T}^d)}
    \le
    C_{\dot\Pi}\|z_t\|_{L^2(\mathbb{T}^d)}
    \qquad\text{for a.e. }t\in[0,1].
    \]
\end{enumerate}
By norm equivalence under the density bounds, the same conclusions hold in
$L^2(\nu)$ up to multiplicative constants depending only on $\lambda,\Lambda$.
\end{restatable}

\begin{proof}
We will start with the Lipschitz bound on the map $t \mapsto \Pi_{\mu_t}z$. Fix $z\in L^2(\mathbb{T}^d;\mathbb R^d)$ and define $u_t \triangleq \Pi_{\mu_t}z$. By first order optimality conditions, $u_t=\nabla \psi_t$ and $\psi_t$ satisfy
\[
\nabla\cdot(\rho_t u_t) = \nabla\cdot(\rho_t z).
\]
Equivalently, for every $\varphi\in H^1(\mathbb{T}^d)$,
\[
\int_{\mathbb{T}^d} \rho_t \langle u_t,\nabla\varphi\rangle\,dx
=
\int_{\mathbb{T}^d} \rho_t \langle z,\nabla\varphi\rangle\,dx.
\]
We first demonstrate the standard $L^2$ bound on the projection. Since $u_t$ is the
$L^2(\mu_t)$-orthogonal projection of $z$ onto the closed subspace
$T_{\mu_t}\mathcal P_2(\mathbb{T}^d)$, we have
\[
\|u_t\|_{L^2(\mu_t)}\le \|z\|_{L^2(\mu_t)}.
\]
Using the density bounds,
\[
\lambda \|u_t\|_{L^2(\mathbb{T}^d)}^2
\le
\|u_t\|_{L^2(\mu_t)}^2
\le
\|z\|_{L^2(\mu_t)}^2
\le
\Lambda \|z\|_{L^2(\mathbb{T}^d)}^2,
\]
hence
\[
\|u_t\|_{L^2(\mathbb{T}^d)}
\le
\sqrt{\frac{\Lambda}{\lambda}}\,\|z\|_{L^2(\mathbb{T}^d)}.
\]
We now prove that the map $t \mapsto \Pi_{\mu_t}z$ is Lipschitz in $L^2(\mathbb{T}^d; \mathbb{R}^d).$ Let $h > 0$ be such that $[t, t+h] \subseteq [0,1]$. Applying the first order conditions at times $t$ and $t+h$ and subtracting gives,
\[\int_{\mathbb{T}^d}\rho_{t+h} \, \langle u_{t+h} - u_t, \nabla \varphi\rangle \, dx = \int_{\mathbb{T}^d}(\rho_{t+h}-\rho_t)\langle z-u_t, \nabla \varphi \rangle\, dx\]
for all $\varphi \in H^1(\mathbb{T}^d).$ Since $u_s=\nabla\psi_s$ for each $s$, define
\[
\eta_h:=\frac{\psi_{t+h}-\psi_t}{h}\in H^1(\mathbb{T}^d),
\qquad
\nabla\eta_h=\frac{u_{t+h}-u_t}{h}.
\]
Taking $\varphi=\eta_h$ in the preceding identity and dividing by $h$ yields
\[
\int_{\mathbb{T}^d}\rho_{t+h}\left\|\frac{u_{t+h}-u_t}{h}\right\|_2^2\,dx
=
\int_{\mathbb{T}^d}\frac{\rho_{t+h}-\rho_t}{h}
\left\langle z-u_t,\frac{u_{t+h}-u_t}{h}\right\rangle\,dx.
\]
Using the lower density bound and applying H\"older's inequality,
\[
\lambda\left\|\frac{u_{t+h}-u_t}{h}\right\|_{L^2(\mathbb{T}^d)}^2
\le
\left\|\frac{\rho_{t+h}-\rho_t}{h}\right\|_{L^\infty(\mathbb{T}^d)}
\|z-u_t\|_{L^2(\mathbb{T}^d)}
\left\|\frac{u_{t+h}-u_t}{h}\right\|_{L^2(\mathbb{T}^d)}.
\]
Therefore,
\[
\left\|\frac{u_{t+h}-u_t}{h}\right\|_{L^2(\mathbb{T}^d)}
\le
\frac{1}{\lambda}
\left\|\frac{\rho_{t+h}-\rho_t}{h}\right\|_{L^\infty(\mathbb{T}^d)}
\|z-u_t\|_{L^2(\mathbb{T}^d)}.
\]
Since $\sup_{t\in [0,1]}\|\partial_t\rho_t\|_{L^\infty(\mathbb{T}^d)} \leq K_\rho$ and $t \mapsto \rho_t$ is absolutely continuous as an $L^\infty(\mathbb{T}^d)$ valued map, we have
\[\|\rho_{t+h} - \rho_t\|_{L^\infty(\mathbb{T}^d)} \leq K_\rho h\]
and thus,
\[
\left\|\frac{u_{t+h}-u_t}{h}\right\|_{L^2(\mathbb{T}^d)}
\le
\frac{K_\rho}{\lambda}\|z-u_t\|_{L^2(\mathbb{T}^d)}.
\]
Applying Minkowski's inequality and the bound on $\|u_t\|_{L^2(\mathbb{T}^d)}$ yields
\[
\|z-u_t\|_{L^2(\mathbb{T}^d)}
\le
\|z\|_{L^2(\mathbb{T}^d)}+\|u_t\|_{L^2(\mathbb{T}^d)}
\le
\left(1+\sqrt{\frac{\Lambda}{\lambda}}\right)\|z\|_{L^2(\mathbb{T}^d)}.
\]
Combining and multiplying through by $h$ therefore results in 
\[
\|u_{t+h}-u_t\|_{L^2(\mathbb{T}^d)}
\le
\frac{K_\rho}{\lambda}
\left(1+\sqrt{\frac{\Lambda}{\lambda}}\right)|h|\,\|z\|_{L^2(\mathbb{T}^d)}.
\]
Thus, statement (1) holds with \[C_{\dot \Pi} = \frac{K_\rho}{\lambda}
\left(1+\sqrt{\frac{\Lambda}{\lambda}}\right).\]
Therefore the map $t \mapsto \Pi_{\mu_t}z$ is Lipschitz as a map from $[0,1]$ to $L^2(\mathbb{T}^d; \mathbb{R}^d)$. Taking the supremum over all $z$ such that $\|z\|_{L^2(\mathbb{T}^d)} = 1$ yields the operator norm bound 
\[
\|\Pi_{\mu_t} - \Pi_{\mu_s}\|_{\mathcal{L}(L^2(\mathbb{T}^d))}
\le
C_{\dot \Pi}|t - s|.
\]
Now let $z_{(\cdot)} \in \text{AC}([0,1]; L^2(\mathbb{T}^d; \mathbb{R}^d))$ and define $y_t \triangleq \Pi_{\mu_t}(z_t).$ Since $z_{(\cdot)}$ is absolutely continuous, $z_t - z_s= \int_{s}^t \partial_r z_r \, dr$. Note that
\[y_t - y_s = (\Pi_{\mu_t} - \Pi_{\mu_s})(z_s) + \Pi_{\mu_t}(z_t-z_s)\]
implies the bound
\[
\|y_t-y_s\|_{L^2(\mathbb{T}^d)}
\le
C_{\dot\Pi}|t-s|\,\|z_s\|_{L^2(\mathbb{T}^d)}
+
\|\Pi_{\mu_t}\|_{\mathcal L(L^2(\mathbb{T}^d))}\int_s^t\|\partial_r z_r\|_{L^2(\mathbb{T}^d)}\,dr.
\]
Since
\[
\|\Pi_{\mu_t}\|_{\mathcal L(L^2(\mathbb{T}^d;\mathbb R^d))}
\le \sqrt{\frac{\Lambda}{\lambda}}
\]
and $z_{(\cdot)}$ is continuous on $[0,1]$, the right-hand side is controlled by
an $L^1$ function of $r$. Therefore $y_{(\cdot)}\in AC([0,1];L^2(\mathbb{T}^d;\mathbb R^d))$. Next define
\[
w_t
\triangleq
y_t-\int_0^t \Pi_{\mu_r}(\partial_r z_r)\,dr.
\]
Since both $y_{(\cdot)}$ and the integral term are absolutely continuous,
$w_{(\cdot)}$ is absolutely continuous as well. Moreover, for $0\le s<t\le 1$,
\begin{align*}
w_t-w_s
&=
\Pi_{\mu_t}z_t-\Pi_{\mu_s}z_s-\int_s^t \Pi_{\mu_r}(\partial_r z_r)\,dr \\
&=
(\Pi_{\mu_t}-\Pi_{\mu_s})z_s
+
\Pi_{\mu_t}\!\left(\int_s^t \partial_r z_r\,dr\right)
-
\int_s^t \Pi_{\mu_r}(\partial_r z_r)\,dr \\
&=
(\Pi_{\mu_t}-\Pi_{\mu_s})z_s
+
\int_s^t(\Pi_{\mu_t}-\Pi_{\mu_r})(\partial_r z_r)\,dr.
\end{align*}
Therefore,
\begin{align*}
\|w_t-w_s\|_{L^2(\mathbb{T}^d)}
&\le
C_{\dot\Pi}|t-s|\,\|z_s\|_{L^2(\mathbb{T}^d)}
+
C_{\dot\Pi}\int_s^t |t-r|\,\|\partial_r z_r\|_{L^2(\mathbb{T}^d)}\,dr.
\end{align*}
Since $w_{(\cdot)}$ is absolutely continuous, it is differentiable for a.e. $t$.
For such $t$, define
\[
(\partial_t\Pi_t)z_t \triangleq \partial_t w_t
=
\partial_t(\Pi_{\mu_t}z_t)-\Pi_{\mu_t}(\partial_t z_t).
\]
It remains to prove the bound on this defect term. Let $t$ be a point at which $w_{(\cdot)}$ is differentiable. For $h$ such that
$[t,t+h]\subset [0,1]$, the identity above gives
\[
w_{t+h}-w_t
=
(\Pi_{\mu_{t+h}}-\Pi_{\mu_t})z_t
+
\int_t^{t+h}(\Pi_{\mu_{t+h}}-\Pi_{\mu_r})(\partial_r z_r)\,dr.
\]
Dividing by $h$ and taking norms,
\[
\left\|\frac{w_{t+h}-w_t}{h}\right\|_{L^2(\mathbb{T}^d)}
\le
C_{\dot\Pi}\|z_t\|_{L^2(\mathbb{T}^d)}
+
C_{\dot\Pi}\frac{1}{|h|}
\int_t^{t+h}|t+h-r|\,\|\partial_r z_r\|_{L^2(\mathbb{T}^d)}\,dr.
\]
Since $|t+h-r|\le |h|$ on the interval of integration,
\[
\frac{1}{|h|}
\int_t^{t+h}|t+h-r|\,\|\partial_r z_r\|_{L^2(\mathbb{T}^d)}\,dr
\le
\int_t^{t+h}\|\partial_r z_r\|_{L^2(\mathbb{T}^d)}\,dr
\to 0
\]
as $h\to 0$. Hence
\[
\|\partial_t w_t\|_{L^2(\mathbb{T}^d)}
\le
C_{\dot\Pi}\|z_t\|_{L^2(\mathbb{T}^d)}
\]
for a.e. $t$, i.e.
\[
\|(\partial_t\Pi_t)z_t\|_{L^2(\mathbb{T}^d)}
\le
C_{\dot\Pi}\|z_t\|_{L^2(\mathbb{T}^d)}.
\]
This proves item (2). The final statement in $L^2(\nu)$ follows from the uniform equivalence between
$L^2(\mathbb{T}^d)$ and $L^2(\nu)$ under the density bounds.
\end{proof}

\end{document}
